\subsection{Batching every proper complex residual}
\label{sec:nested-complex}

We retain the complex scalar circuit, binary labels and shared auxiliary
banks of PR~\#7, and the whole-residual identity and unequal-width row
construction of PR~\#10. The change is to apply that identity to every
nonzero physical residual, instead of only the two large auxiliary classes.
The complete child list requires a new arithmetic-path guard, proved below.
Put
\[
 h=28,\quad v=\binom{28}{3}=3276,\quad m=h^3=21952,\quad
 R=93838,\quad N=v^3,\quad B=v^2(R+h+1).
\]
Here $R$ counts side roles and the additional $h+1$ roles are the unchanged
point and total centers. Thus
\[
 B=1007397164592,\qquad W=2v^2(v+R+h+1)=2085111546336,
\]
\[
 L=3v^2h(h+1),\qquad D=2N-2L=18030055680,\qquad
 s=Wm-D=45772350635112192.
\]

\paragraph{All physical edge ranks.}
Run the compiled forward and inverse invocation schedules with local source
labels of dimension $1$ on $X$ and $0$ on $Y$ and the auxiliary roles, and
local sink labels of dimension $h$ on $X$ and the auxiliary roles and $h-1$
on $Y$. At a scalar gate, count the change in dimension on every touched
physical role, including a nonpivot input whose value is unchanged.
The two directions give the same local histogram $c_r$. Its nonzero part is
\[
\begin{array}{r|r@{\qquad}r|r@{\qquad}r|r@{\qquad}r|r}
 r&c_r&r&c_r&r&c_r&r&c_r\\\hline
1 & 93912 & 8 & 10912 & 15 & 4536 & 22 & 5901\\
2 & 23891 & 9 & 5656 & 16 & 7803 & 23 & 8400\\
3 & 21224 & 10 & 8456 & 17 & 4228 & 24 & 11550\\
4 & 31667 & 11 & 4284 & 18 & 8232 & 25 & 1512\\
5 & 13552 & 12 & 7598 & 19 & 4872 & 26 & 5180\\
6 & 9128 & 13 & 4872 & 20 & 7482 & 27 & 6552\\
7 & 9240 & 14 & 7448 & 21 & 7224 & 28 & 87\\
\end{array}
\]
There are additionally $727056$ zero-rank physical edges per invocation.
The exact sum is
\[
 \sum_{r=1}^{28}r c_r
  =(2v+R+h+1)h-2v+2h(h+1)=2806804.
\]
Every decreasing edge is one of the $h+1$ central returns on a single
chronological frontier, and has loss $h$.

At tensor stage $j$, set $a=h^{j-1}$. The common local frame of dimension
$u$ is embedded as a frame of dimension $(a-1)h+u$. Auxiliary frames run
from $D_0$ of dimension $(a-1)h$ to $D_1$ of dimension $ah$. Physical data
enters at dimensions $a$ and $a-1$, while its first relevant gates have
dimensions $(a-1)h+1$ and $(a-1)h$. Both entrance ranks are therefore
$(a-1)(h-1)$. This calculation also applies to the inverse middle stage:
its first source copy is at the low frame, and its first target injection
is at the line frame. Their local entrance ranks are zero, so adding the
following entrance classes does not count any local edge twice.

The global histogram $H_r$ consists of $3v^2$ copies of $c_r$ and the five
additional classes
\[
\begin{array}{l|r|r}
 \text{class}&\text{copies}&\text{rank}\\\hline
 \text{shared stage-one/stage-three auxiliary join}&B&m-2h=21896\\
 \text{stage-two auxiliary sink}&B&m-h^2=21168\\
 \text{stage-two auxiliary entrance}&B&h^2-h=756\\
 \text{stage-two data entrance}&2N&(h-1)^2=729\\
 \text{stage-three data entrance}&2N&(h^2-1)(h-1)=21141.
\end{array}
\]
There is no extra entrance for a stage-one auxiliary role, because $D_0=0$,
and no extra terminal edge for a stage-three auxiliary role, because
$D_1$ is full. The shared bank has one join in place of a separate
stage-one terminal edge and a stage-three entrance. A stage-two bank has
one entrance and one terminal edge. The local data sinks are exactly the
physical data frames entering the next stage. Thus this is a complete
physical-edge list, not merely an equality of total ranks. It satisfies
\[
 \sum_{r\ge1}rH_r=s,
 \qquad 1\le r\le M:=21896<m\quad\text{whenever }H_r>0.
\]
In particular there are no rank-$m$ exceptions.
The reproducible checker is
\texttt{research/nested-source/complex\_all.py}.

\paragraph{Residual bases, both directions, and endpoints.}
The inherited exact binary checker verifies all local nondegenerate labels
and nonalternating nonzero residuals, and follows each compiled role in
both directions. It also verifies that the dimensions used in the histogram
are their exact binary ranks. Tensoring the checked orthonormal residual
bases with the norm-one past and future line factors preserves
orthonormality.

For completeness the external classes have the same property. Write the
past space as $A=B_0\mathbin\perp P$, with $P$ its norm-one tensor-indicator
line. At stages two and three, $B_0=P^\perp$ is nondegenerate and contains
a coordinate unit outside the support of that indicator. It therefore
has an orthonormal basis. The data entrance residual is
\[
 B_0\otimes t_T^\perp\otimes Q,
\]
where $Q$ is a norm-one future line, and $t_T^\perp$ is nondegenerate and
contains coordinate units outside $T$. The stage-two auxiliary entrance
has residual $B_0\otimes F\otimes Q$, and its sink has residual
$F\otimes F\otimes Q^\perp$. These products all have orthonormal bases.
The inverse-stage entrances are at the same physical low or line frames
just identified, so this argument covers that direction as well.

For the shared-bank join use the retained point-partner triple permutation
$\pi$. Its even intersections give
\[
 E=F\otimes\langle t_A\rangle\otimes\langle t_B\rangle
 \quad\subset\quad
 H=\langle t_B\otimes t_{\pi(A)}\rangle^\perp\otimes F.
\]
The nondegenerate residual $H\cap E^\perp$ contains
$e_i\otimes e_j\otimes e_k$ with $i,k\notin B$, a vector of norm one.
It consequently has an orthonormal basis. This also explains the rank
$m-2h$ without appealing only to dimensions.
The scalar circuit, arbitrary auxiliary restoration, terminal frames and
fourth-root endpoint corrections have not changed.

\paragraph{The enlarged time moment.}
On an edge of rank $r$, keep its existing binary basis maps, placing the
$r$ orthonormal residual slots first. At an integer parent width
$e=mf+t$, let $B_-$ index the negative residual directions. The identity
$C^{-1}=-iZCZ$ gives
\[
 \bigotimes_{j=1}^{r}(C^{\varepsilon_j})^{\otimes f}
   =(-i)^{f|B_-|}Z_{B_-}\,C^{\otimes rf}\,Z_{B_-}.
\]
One child of width $rf$, two address-sign passes and one fourth-root
phase thus replace all directional children on the edge. This is the
whole-residual construction of PR~\#10 applied to the complete finite list.
All children contract strictly, since $r<m$. Its row reservation and
integer-width construction apply unchanged, including the individually
processed fewer than $m$ remainder axes and volume factor $1/W$ per child.
No simultaneous change of binary basis on different edges is needed: each
edge retains its own original frame conjugations.

The resulting normalized time moment is
\[
 \Psi_{\rm c}(1-a_{\rm c})
   =\sum_{r\ge1}\frac{H_r r}{Wm}
          \exp\!\left(a_{\rm c}\log\frac mr\right),
 \qquad \sum_{r\ge1}\frac{H_r r}{Wm}=1-\frac D{Wm}.
\]
The certificate encloses each logarithm by the positive rational series
for $\log x$ after extracting a power of two. With upper bounds $\ell_r$,
\[
 \Psi_{\rm c}(1-a_{\rm c})
 <\sum_{r\ge1}\frac{H_r r}{Wm(1-a_{\rm c}\ell_r)}<1
 \qquad\left(a_{\rm c}=\frac4{10^6}\right).
\]
The last comparison uses exact fractions; its gap exceeds
$1.8\cdot10^{-8}$. The certificate also supports $a_{\rm c}=419/10^8$,
but the displayed smaller saving leaves more slack.
The unchanged scalar map is independently rechecked on all $h=28$ source
coefficients by \texttt{complex\_controls.py}, together with exact mixed-sign
Gaussian-dyadic and arbitrary-dirty auxiliary controls.

\paragraph{A path guard for all residual children.}
\label{sec:nested-complex-guard}
The common-gate potential argument of PR~\#10 gives
\[
 \sum_{e\in P} r_e
 =\dim U_{\rm end}-\dim U_{\rm start}
  +2\sum_{e\in P\text{ decreasing}}
       (\dim U_e^- -\dim U_e^+)
 \le q:=m+6h=22120
\]
on every scalar dependency path. In this identity all dimensions are
\emph{global} tensor-frame dimensions; one does not restart the potential
at a stage boundary. Every data entrance, stage-two auxiliary entrance or
sink, and shared-bank join is increasing, so each is already charged once
in the telescoping endpoint difference. Only central returns decrease.
A path meets at most one such frontier per tensor stage, because the
invocations within a stage have disjoint data and auxiliary wires.
Sharing banks connects stages one and three chronologically and introduces
no additional return. Thus the total decrease on a path is at most $3h$.

The coefficient-path projection remains valid after whole-residual
batching. A child acts within one parent role stream, binary basis maps
only permute its addresses, and row splitting and parking create no
arithmetic dependencies between parent roles. A path may change parent
roles only at a scalar gate with its common frame. Hence its child widths
are $r_e f$ with $r_e\le M$ and $\sum_e r_e\le q$.

For $\rho>1$, convexity and $M<q<2M$ give, for every such list,
\[
 \sum_e\left(\frac{r_e}{m}\right)^\rho
 \le\left(\frac M m\right)^\rho
       +\left(\frac{q-M}{m}\right)^\rho.
\]
Indeed moving mass from a smaller nonzero entry into a larger one, until
the latter reaches $M$, can only increase the sum. The maximal relaxed
list is therefore $(M,q-M)$; this relaxation need not occur as an actual
circuit path.
Take $\rho=3/2$, $t=(m-M)/m=1/392$ and $y=(q-M)/m=1/98$. Taylor's formula
with its second derivative bounded on $[0,t]$ gives
\[
 (1-t)^{3/2}\le1-\frac32t+\frac38\frac{t^2}{1-t},
 \qquad y^{3/2}<\frac y9.
\]
Their sum is exactly bounded by
\[
 \frac{11005895}{11035584}<\frac{999}{1000}=:\overline\theta<1.
\]
Integer floors only shrink children. This proves the required path moment
for \emph{every} residual class, including the intermediate-size entrance
edges; it does not assume that all edges outside the old two bulk classes
remain singleton children.

Set
\[
 E=64(W+m+1)^3,\qquad C_{\rm dep}=1000(E+16m+1).
\]
The exact inequality
\[
 36W^3+4s+4W+8m+4<E
\]
covers the scalar operations, sign wrappers, and remainder-axis kernels.
There are at most $s$ nonzero edges, so the existing $4s$ allowance also
covers wrappers on the enlarged list. A stopping leaf reached from an
internal call has width at least $d^\beta/(2m)$. Since
$8(2m)^{\rho-1}\le16m$ and
$(1-\overline\theta)C_{\rm dep}\ge E$, the same strong induction as in
PR~\#10 gives
\[
 A(e)\le C_{\rm dep}\,d^{-\beta(\rho-1)}e^\rho.
\]
With the retained completed-layer allowances, take
\[
 C_1=\rho-(\rho-1)\beta+\zeta,
 \qquad
 C_0=\left\lceil128m(1+1/\zeta)C_{\rm dep}\right\rceil.
\]
For $\beta=1/1000$ and $\zeta=1/10000$ this yields
\[
 C_1=\frac{3749}{2500}=1.4996.
\]
Thus the expanded complex recurrence has the same finite-alphabet phase
interface and a proved precision guard, with its stronger time exponent.
The integer-multiplication claim requires the separate bit interface and
assembly certificate; a complex-layer saving alone is not such a claim.
